%\begin{}
\paragraph{ }
A processor is a specialized digital circuit designed to execute instructions, perform mathematical and logical calculations, and control data flow within a computer system. Over the past decades, microprocessors have not only become a defining trend in the technology industry but have fundamentally reshaped it. As a result, studying and advancing processor design remains a priority for both industry and academia.
\paragraph{}
RISC-V (Reduced Instruction Set Computer - Five) is an open-standard instruction set architecture (ISA) developed in 2010 at UC Berkeley by Krste Asanović, Andrew Waterman, and Yunsup Lee under the supervision of David Patterson. What makes it unique is its open-standard nature, allowing users to adopt it freely and even modify it, fostering rapid adoption across both academic and commercial domains. Additionally, its adherence to the RISC philosophy makes it not only easier to design and implement but also inherently energy efficient.
\paragraph{}
This project presents BigDCore (Big Dream Core), a pipelined 32-bit processor implementing the full RV32I Base Integer ISA, described in VHDL and targeting simulation-based verification. RV32I was chosen as the target ISA for its completeness as a base integer instruction set while remaining tractable for a single academic project. Testing is conducted using the industry-standard simulation tool ModelSim, and future work includes the privileged architecture to enable an operating system to run on the processor.
